\documentclass[10pt]{article}
\usepackage[margin=0.95in]{geometry}
\usepackage{amsmath,amssymb}
\usepackage[T1]{fontenc}
\usepackage{microtype}
\setlength{\parskip}{2pt}
\usepackage{multicol}
\title{\vspace{-1.2em}\large\textbf{Readings memo: equivalence scales and sequential fertility}\vspace{-0.6em}}
\author{\normalsize Fertility--housing project}
\date{\normalsize July 18, 2026 --- readings only; no model change is proposed here}
\begin{document}
\maketitle
\vspace{-1.6em}

\paragraph{Why these readings.}
The model's income risk is capped by feasibility: every household must afford
$\bar c(n) + r\bar h(n)$ each period, the childless bundle is roughly a third
of mean income, and no measured safety-net guarantee for childless adults
comes close to that. The July~18 frontier computation makes the bind exact:
at the calibrated persistence the model is infeasible above
$\sigma_z\!\approx\!0.09$ without intervention; the measured floor buys no
room at the current intercept; lowering $\bar c_0$ to $0.10$ of mean income
plus the floor reaches $\sigma_z=0.20$, where the estate tail nearly closes
but young liquid wealth overshoots its target seven-fold. Two literatures
speak to the two structural exits. Equivalence scales replace subtraction
with division in preferences, so no income realization is ever infeasible.
Sequential-fertility models make the \emph{timing} of births a priced choice,
which changes how fertility absorbs risk and housing costs and which moments
identify it.

\paragraph{1. Equivalence scales.}
An equivalence scale is a divisor $m(n)$ that converts household spending
into a per-adult-equivalent standard of living. The expert scales are
conventions: the OECD ``Oxford'' scale assigns $1$ to the first adult, $0.7$
to further adults, $0.5$ to a child; the OECD-modified scale (Hagenaars, de
Vos and Zaidi 1994) uses $0.5$ and $0.3$; the square-root scale divides by
$\sqrt{\text{household size}}$. The measurement tradition estimates $m$ from
behavior: Engel's method reads adult welfare off the food budget share,
Rothbarth's off spending on adult goods, and the two disagree wildly ---
Deaton and Muellbauer (1986) report that under the most extreme assumptions
Engel makes children about four times as expensive as Rothbarth. Barten
(1964) scales each commodity separately, so children change relative shadow
prices, not just a single deflator. The identification limit is the
discipline's oldest warning: Pollak and Wales (1979) showed that
welfare-relevant scale \emph{levels} are not identified from demand data
alone, and Blundell and Lewbel (1991) sharpened this to ``only changes, not
levels.'' What survives is either an untestable restriction --- Lewbel's
(1989) base independence, under which the scale is a utility-free function of
demographics recoverable from shape-invariant demands --- or a change of
question: Browning, Chiappori and Lewbel (2013) estimate ``indifference
scales'' inside a collective model and find singles need one half to three
quarters of a couple's outlay, couples saving about a third through shared
consumption. Quantitative life-cycle work sidesteps estimation and imposes a
scale: Fern\'andez-Villaverde and Krueger (2007) deflate consumption by the
mean of six published scales; Attanasio, Banks, Meghir and Weber (1999) put
demographics in the Euler equation as taste shifters; and --- the bridge case
for us --- Scholz, Seshadri and Khitatrakun (2006) use
$n_j\,U(c_j/n_j)$ with $n_j=(A_j+0.7K_j)^{0.7}$ \emph{and} scale their
means-tested floor by the same $n_j$, i.e.\ scale-in-utility and floor-in-institutions
coexist in a published JPE quantitative model.

For our problem the arithmetic is the point. With $u(c/m(n))$ and no
translation term, utility is defined for every $c>0$: however low the income
draw, a feasible choice exists, the empty-budget-set region vanishes
identically, and $\sigma_z$ becomes an estimable parameter rather than a
feasibility bound. The honest cost is that our nonhomotheticity lives in the
intercepts: committed child expenditures make children and housing
necessities, which is where the fertility--income and tenure gradients come
from. Under pure scaling, child costs turn proportional to the standard of
living --- a Becker--Lewis flavored substitution, but a different one, and
the whole fertility block re-identifies. A modern middle path exists:
nonhomothetic CES (Comin, Lashkari and Mestieri 2021) has \emph{no}
subtraction term yet generates non-vanishing necessity/luxury structure
(their relative Engel elasticities are constants, where Stone--Geary income
effects ``level off quickly''). Two negative results from this pass: we
found no structural model with a housing-specific equivalence scale, and no
paper stating the no-empty-budget-set property explicitly --- the property is
elementary, but we would be asserting it ourselves.

\paragraph{2. Sequential fertility.}
The dynamic-programming lineage treats fertility as a repeated discrete
choice, not a completed quantity. Wolpin (1984) is the canonical estimated
model: a per-period birth decision from age 15 to 45 under child mortality.
Hotz and Miller (1988) add imperfect conception control, with child-care
costs shaping spacing; Arroyo and Zhang (1997) survey the class. Adda,
Dustmann and Stevens (2017) put sequential births inside occupation, labor
supply and savings choices and price the career cost of children at 35.3\%
of lifetime income, three quarters of it from employment interruptions. Two
papers sit closest to us. Sommer (2016) embeds a binary per-period birth
choice in a Bewley model where infertility risk rises exponentially with age
(reaching one at 45) and finds that realistic increases in uninsurable
earnings risk postpone births and reduce their number --- the tempo and
quantum margins moving together. Ejrn\ae s and J\o rgensen (2020) estimate
sequential contraceptive effort with income risk and find abortion access
functions as an insurance margin. Doepke and Kindermann (2019) make each
birth require both partners' agreement, a veto structure that matters for
policy targeting; Baudin, de la Croix and Gobbi (2015) split childlessness
into poverty-driven (2.5\% of women) and opportunity-cost-driven (8.1\%)
components. De la Croix and Pommeret (2021) make the option value explicit:
postponing a birth is delaying an irreversible investment, and enough
uncertainty pushes first births past a threshold --- against an
age-declining fecundity schedule calibrated to demographic data. The
demographers' vocabulary is Bongaarts and Feeney (1998): quantum is the
lifetime number, tempo the mean age, and observed TFR confounds them. The
reduced-form housing evidence is timing-shaped too: Dettling and Kearney
(2014) find house-price increases raise owner births ($+5\%$ per \$10{,}000)
and depress renter births ($-2.4\%$); Lovenheim and Mumford (2013) find
\$100{,}000 of housing wealth raises the probability of a birth by 16--18\%
with no renter effect.

Read against our architecture: we already have period-by-period births in a
fecund window with the $(n,s)$ child states, so ``sequential'' is not the
missing piece --- \emph{pricing time} is. Without age-declining fecundity,
waiting is costless, so the option value of postponement is mispriced and
risk responses load entirely on the number of children (our $\sigma_z=0.20$
probe raises childlessness from $0.19$ to $0.28$ at fixed parameters); with
a biological clock, part of that response becomes delayed first births and
compressed spacing, which is what the data's timing moments measure. A
fecundity schedule is two externally calibrated parameters, and it converts
unused moments --- age at first birth, spacing, birth hazards by age and by
number of children already born --- into identifying variation. It would
also let the second-birth housing-response target be what its name says,
rather than housing-demand discipline.

\paragraph{What this buys, jointly.}
The two readings answer different halves of one problem. Scales (or
nonhomothetic CES) remove the feasibility cap on $\sigma_z$ at the cost of
re-identifying the fertility block; a fecundity-priced sequential structure
supplies precisely the new timing moments that block would need, and gives
risk a tempo margin our current model lacks. They are complements, each
separable, and both are redesigns --- nothing here modifies the model today.
The near-term income-process question (tail versus young liquid wealth)
stands unchanged in the meantime.

\begin{footnotesize}
\begin{multicols}{2}
\noindent
Adda, Dustmann, Stevens (2017), \emph{JPE} 125(2).\\
Arroyo, Zhang (1997), \emph{J.\ Population Economics} 10(1).\\
Attanasio, Banks, Meghir, Weber (1999), \emph{JBES} 17(1).\\
Barten (1964), in \emph{Econometric Analysis for National Economic Planning}.\\
Baudin, de la Croix, Gobbi (2015), \emph{AER} 105(6).\\
Blundell, Lewbel (1991), \emph{J.\ Econometrics} 50.\\
Bongaarts, Feeney (1998), \emph{PDR} 24(2).\\
Browning (1992), \emph{JEL} 30(3).\\
Browning, Chiappori, Lewbel (2013), \emph{ReStud} 80(4).\\
Comin, Lashkari, Mestieri (2021), \emph{Econometrica} 89(1).\\
Deaton, Muellbauer (1986), \emph{JPE} 94(4).\\
de la Croix, Pommeret (2021), \emph{JET} 193.\\
Dettling, Kearney (2014), \emph{J.\ Public Economics} 110.\\
Doepke, Kindermann (2019), \emph{AER} 109(9).\\
Ejrn\ae s, J\o rgensen (2020), \emph{J.\ Applied Econometrics} 35(5).\\
Fern\'andez-Villaverde, Krueger (2007), \emph{ReStat} 89(3).\\
Hagenaars, de Vos, Zaidi (1994), \emph{Poverty Statistics in the Late 1980s}.\\
Hotz, Miller (1988), \emph{Econometrica} 56(1).\\
Lewbel (1989), \emph{J.\ Public Economics} 39.\\
Lovenheim, Mumford (2013), \emph{ReStat} 95(2).\\
Pollak, Wales (1979), \emph{AER} 69(2).\\
Scholz, Seshadri, Khitatrakun (2006), \emph{JPE} 114(4).\\
Sommer (2016), \emph{JME} 83.\\
Wolpin (1984), \emph{JPE} 92(5).
\end{multicols}
\end{footnotesize}

\end{document}
